\documentclass[10pt,a4paper]{amsart}
\usepackage{amsmath,amssymb,amsthm,graphicx}
% Retained minimal preamble for the Stacks Project chapter more-morphisms.tex (More on Morphisms).
% Verbatim excerpt of the chapter preamble: Xy-pic setup (native preamble lines 14-18) and the native
% theorem declarations with the mathematical macros (native preamble lines 154-217). Copied mechanically;
% nothing added, removed or redefined. The wrapper class and the amsmath/amssymb/amsthm/graphicx base
% packages are supplied by the assembler, so the class/externaldocument part of the native preamble is omitted.
% For commutative diagrams we use Xy-pic
\usepackage[all]{xy}

% We use 2cell for 2-commutative diagrams.
\xyoption{2cell}
% Theorem environments.
%
\theoremstyle{plain}
\newtheorem{theorem}[subsection]{Theorem}
\newtheorem{proposition}[subsection]{Proposition}
\newtheorem{lemma}[subsection]{Lemma}

\theoremstyle{definition}
\newtheorem{definition}[subsection]{Definition}
\newtheorem{example}[subsection]{Example}
\newtheorem{exercise}[subsection]{Exercise}
\newtheorem{situation}[subsection]{Situation}

\theoremstyle{remark}
\newtheorem{remark}[subsection]{Remark}
\newtheorem{remarks}[subsection]{Remarks}

\numberwithin{equation}{subsection}

% Macros
%
\def\lim{\mathop{\mathrm{lim}}\nolimits}
\def\colim{\mathop{\mathrm{colim}}\nolimits}
\def\Spec{\mathop{\mathrm{Spec}}}
\def\Hom{\mathop{\mathrm{Hom}}\nolimits}
\def\Ext{\mathop{\mathrm{Ext}}\nolimits}
\def\SheafHom{\mathop{\mathcal{H}\!\mathit{om}}\nolimits}
\def\SheafExt{\mathop{\mathcal{E}\!\mathit{xt}}\nolimits}
\def\Sch{\mathit{Sch}}
\def\Mor{\mathop{\mathrm{Mor}}\nolimits}
\def\Ob{\mathop{\mathrm{Ob}}\nolimits}
\def\Sh{\mathop{\mathit{Sh}}\nolimits}
\def\NL{\mathop{N\!L}\nolimits}
\def\CH{\mathop{\mathrm{CH}}\nolimits}
\def\proetale{{pro\text{-}\acute{e}tale}}
\def\etale{{\acute{e}tale}}
\def\QCoh{\mathit{QCoh}}
\def\Ker{\mathop{\mathrm{Ker}}}
\def\Im{\mathop{\mathrm{Im}}}
\def\Coker{\mathop{\mathrm{Coker}}}
\def\Coim{\mathop{\mathrm{Coim}}}

% Boxtimes
%
\DeclareMathSymbol{\boxtimes}{\mathbin}{AMSa}{"02}

%
% Macros for moduli stacks/spaces
%
\def\QCohstack{\mathcal{QC}\!\mathit{oh}}
\def\Cohstack{\mathcal{C}\!\mathit{oh}}
\def\Spacesstack{\mathcal{S}\!\mathit{paces}}
\def\Quotfunctor{\mathrm{Quot}}
\def\Hilbfunctor{\mathrm{Hilb}}
\def\Curvesstack{\mathcal{C}\!\mathit{urves}}
\def\Polarizedstack{\mathcal{P}\!\mathit{olarized}}
\def\Complexesstack{\mathcal{C}\!\mathit{omplexes}}
% \Pic is the operator that assigns to X its picard group, usage \Pic(X)
% \Picardstack_{X/B} denotes the Picard stack of X over B
% \Picardfunctor_{X/B} denotes the Picard functor of X over B
\def\Pic{\mathop{\mathrm{Pic}}\nolimits}
\def\Picardstack{\mathcal{P}\!\mathit{ic}}
\def\Picardfunctor{\mathrm{Pic}}
\def\Deformationcategory{\mathcal{D}\!\mathit{ef}}
\usepackage{cleveref}
\crefname{originallemma}{original lemma}{original lemmas}
\Crefname{originallemma}{Original lemma}{Original lemmas}
\makeatletter
\newlabel{lemma-ind-quasi-affine-alternative-definition}{{0F1U}{source}{}{sourcepointer-1}{}}
\newlabel{lemma-ind-quasi-affine-alternative-definition@cref}{{[originallemma][0][]0F1U}{[1][0][]source}{}{}{}}
\makeatother
\crefname{originallemma}{original lemma}{original lemmas}
\Crefname{originallemma}{Original lemma}{Original lemmas}
\makeatletter
\newlabel{lemma-ind-quasi-affine-composition}{{0F1V}{source}{}{sourcepointer-2}{}}
\newlabel{lemma-ind-quasi-affine-composition@cref}{{[originallemma][0][]0F1V}{[1][0][]source}{}{}{}}
\makeatother
\crefname{originallemma}{original lemma}{original lemmas}
\Crefname{originallemma}{Original lemma}{Original lemmas}
\makeatletter
\newlabel{lemma-ind-quasi-affine-examples}{{0F1W}{source}{}{sourcepointer-3}{}}
\newlabel{lemma-ind-quasi-affine-examples@cref}{{[originallemma][0][]0F1W}{[1][0][]source}{}{}{}}
\makeatother
\crefname{originallemma}{original lemma}{original lemmas}
\Crefname{originallemma}{Original lemma}{Original lemmas}
\makeatletter
\newlabel{lemma-ind-quasi-affine-permanence}{{0F1X}{source}{}{sourcepointer-4}{}}
\newlabel{lemma-ind-quasi-affine-permanence@cref}{{[originallemma][0][]0F1X}{[1][0][]source}{}{}{}}
\makeatother
\title{Being ind-quasi-affine is stable under base change}
\author{The Stacks Project Authors}
\date{stacks-project:more-morphisms.tex:local-snapshot-20261011; GNU FDL 1.2 or later}
\begin{document}
\maketitle
\noindent\textit{Editorial note.} Counted span is the author's own proof environment, source lines 19639-19652, immediately following the statement at 19634-19637.\par\medskip
\noindent\textit{Editorial note.} The author's proof is a reduction through four cited chapter lemmas; all four are navigation pointers into the same chapter and carry no part of the statement being counted.\par\medskip
\noindent\textit{Editorial note.} The statement itself is short but is not restated anywhere else in this chapter or in its published ranges.\par\medskip
\noindent\textit{Allowed standard prerequisites (editorial).}\begin{itemize}
\item Alternative definition of ind-quasi-affine. chapter lemma lemma-ind-quasi-affine-alternative-definition, cited at proof line 19642. Lets the property be checked Zariski locally on the target, so the base may be assumed affine.
\item Composition of ind-quasi-affine morphisms. chapter lemma lemma-ind-quasi-affine-composition, cited at proof line 19647. Combines the affine base change X x\_Y Z -> X with the given morphism X -> Y.
\item Ind-quasi-affine examples, including affine morphisms. chapter lemma lemma-ind-quasi-affine-examples, cited at proof line 19648. Supplies that the affine morphism X x\_Y Z -> X is ind-quasi-affine.
\item Permanence of ind-quasi-affineness. chapter lemma lemma-ind-quasi-affine-permanence, cited at proof line 19650. Transfers the property from the composite X x\_Y Z -> Y to the base change X x\_Y Z -> Z.
\end{itemize}\medskip
\ifdefined\hypertarget\hypertarget{sourcepointer-1}{}\fi
\noindent\textit{Source reference.} Original lemma 0F1U: sources/stacks-more-morphisms/pinned-original-source.tex, source lines 19538--19538. Navigation pointer only: the counted proof cites this statement by \textbackslash{}ref as a standing input; the deduction the counted proof has to make is carried out inside the counted span, and no part of that deduction is delegated to this pointer.\par\medskip
\ifdefined\hypertarget\hypertarget{sourcepointer-2}{}\fi
\noindent\textit{Source reference.} Original lemma 0F1V: sources/stacks-more-morphisms/pinned-original-source.tex, source lines 19585--19585. Navigation pointer only: the counted proof cites this statement by \textbackslash{}ref as a standing input; the deduction the counted proof has to make is carried out inside the counted span, and no part of that deduction is delegated to this pointer.\par\medskip
\ifdefined\hypertarget\hypertarget{sourcepointer-3}{}\fi
\noindent\textit{Source reference.} Original lemma 0F1W: sources/stacks-more-morphisms/pinned-original-source.tex, source lines 19606--19606. Navigation pointer only: the counted proof cites this statement by \textbackslash{}ref as a standing input; the deduction the counted proof has to make is carried out inside the counted span, and no part of that deduction is delegated to this pointer.\par\medskip
\ifdefined\hypertarget\hypertarget{sourcepointer-4}{}\fi
\noindent\textit{Source reference.} Original lemma 0F1X: sources/stacks-more-morphisms/pinned-original-source.tex, source lines 19621--19621. Navigation pointer only: the counted proof cites this statement by \textbackslash{}ref as a standing input; the deduction the counted proof has to make is carried out inside the counted span, and no part of that deduction is delegated to this pointer.\par\medskip
\begin{lemma}
\label{lemma-base-change-ind-quasi-affine}
The property of being ind-quasi-affine is stable under base change.
\end{lemma}

\begin{proof}
Let $f : X \to Y$ be an ind-quasi-affine morphism.
For checking that every base change of $f$ is ind-quasi-affine, by
Lemma \ref{lemma-ind-quasi-affine-alternative-definition}, we may work
Zariski locally on $Y$, so we assume that $Y$ is affine.
Furthermore, we may also assume that in the base change morphism
$Z \to Y$ the scheme $Z$ is affine, too. The base change
$X \times_Y Z \to X$ is an affine morphism, so, by
Lemmas \ref{lemma-ind-quasi-affine-composition} and
\ref{lemma-ind-quasi-affine-examples},
the map $X \times_Y Z \to Y$ is ind-quasi-affine. Then, by
Lemma \ref{lemma-ind-quasi-affine-permanence}, the
base change $X \times_Y Z \to Z$ is ind-quasi-affine, as desired.
\end{proof}

\end{document}
